% Generated draft: requires mathematical and source review.
\documentclass[11pt]{article}
\usepackage[margin=25mm]{geometry}
\usepackage{fontspec}
\setmainfont{Latin Modern Roman}
\usepackage{amsmath,amssymb,mathtools,unicode-math}
\setmathfont{Latin Modern Math}
\usepackage{enumitem,needspace,xurl,hyperref}
\setlength{\parindent}{0pt}
\setlength{\parskip}{5pt}
\setlength{\emergencystretch}{3em}
\newcommand{\N}{\mathbb{N}}
\newcommand{\Z}{\mathbb{Z}}
\newcommand{\Q}{\mathbb{Q}}
\newcommand{\R}{\mathbb{R}}
\newcommand{\C}{\mathbb{C}}
\newcommand{\F}{\mathbb{F}}
\newcommand{\A}{\mathbb{A}}
\newcommand{\PP}{\mathbb{P}}
\newcommand{\E}{\mathbb{E}}
\newcommand{\eps}{\varepsilon}
\newcommand{\dd}{\,\mathrm{d}}
\begin{document}
\section*{UnsolvedMath ID: 20001861}
\textbf{Source problem code:} catalogue entry
\subsection{Definitions and mathematical statement}
The source attributes to Whiteley a conjecture that the configuration space of a symmetric bar-joint framework, after fixing suitable reference-frame data, has the symmetry of its most symmetric realization. It does not define the symmetry category, most symmetric, frame fixing or faithfulness. The supplied verification corrects the extracted item number from 2 to 12; trailing ``3 1'' and the hyphen in ``configu-ration'' are layout artifacts.

For a precise intrinsic reading, fix a finite graph \(G=(V,E)\), positive bar lengths \(\ell\), and
\[
 \mathcal R_\ell=\{q\in(\R^d)^V:|q_i-q_j|=\ell_{ij}\text{ for }ij\in E\}.
\]
Let \(\operatorname{Aut}(G,\ell)\) preserve graph and length labels, acting by relabeling. The aid reports that this action is isometric on the unquotiented space and descends to both \(\mathcal R_\ell/\operatorname{SE}(d)\) and \(\mathcal R_\ell/\operatorname{E}(d)\), with their quotient metrics. Whether collisions or self-crossings are allowed requires a convention; the displayed realization equations impose only bars.

For the fully Euclidean-reduced space \(\mathcal M=\mathcal R_\ell/\operatorname{E}(d)\) and a full-dimensional realization \(q\), the projected spatial symmetry group is
\[
 H_q=\{\sigma\in\operatorname{Aut}(G,\ell):
 \exists T\in\operatorname{E}(d),\ Tq_i=q_{\sigma(i)}\text{ for every }i\}
 =\operatorname{Stab}([q]).
\]
Full dimension makes the spatial isometry determined by its action on the realization; vertex coincidences would need separate care if identifying this projected group with an unlabelled spatial symmetry group.

Let \(C\) be the connected component through \([q]\), and put
\[
 K_C=\{\sigma\in\operatorname{Aut}(G,\ell):
 \sigma z=z\text{ for every }z\in C\}.
\]
Since \(H_q\) fixes \([q]\), it preserves \(C\). The aid's exact effective relabeling symmetry supplied on that component is
\[
 H_q/(H_q\cap K_C).
\]
It is faithful precisely when \(H_q\cap K_C=1\). This proves a possibly nonfaithful inherited-action interpretation, not equality with the entire intrinsic isometry group or a classification of all most-symmetric realizations.

The aid reports two elementary obstructions to stronger readings. An equilateral triangular framework has fully Euclidean-reduced configuration space consisting of one point, so its nontrivial triangle symmetry group acts trivially. Thus faithfulness is not automatic. For a pinned unit four-bar \(C_4\), a regular parallelogram branch can be written
\[
 q_1=(0,0),\quad q_2=(1,0),\quad
 q_3=(1+\cos\theta,\sin\theta),\quad
 q_4=(\cos\theta,\sin\theta),\quad0<\theta<\pi.
\]
Cyclic relabeling does not directly preserve these pins. After the configuration-dependent Euclidean renormalization restoring them, the reported parameter action is \(\theta\mapsto\pi-\theta\). A fixed coordinate slice therefore requires compensated actions, rather than assuming every relabeling preserves the chosen frame literally. These examples distinguish interpretations of the conjecture; they do not establish a unique intended formulation from the abbreviated source.
\subsection{Short English statement}
Whiteley's symmetry conjecture needs a choice of quotient, frame and faithfulness. The aid gives the exact relabeling kernel quotient on a moduli component. A triangle shows inherited symmetry can be trivial, and a pinned four-bar shows that frame restoration may depend on the configuration.
\subsection{Sources}
\begin{itemize}
\item Local UnsolvedMath registry, permanent ID \(20001861\), code catalogue entry; canonical \texttt{aim-geometry-notes.json}, \texttt{notes[198]}.
\item AIM, ``Configuration spaces of linkages,'' conjecture attributed to Whiteley, \url{https://aimath.org/pastworkshops/linkagesproblems.pdf}. The supplied verification identifies original item12 rather than extracted item2. No PDF inspection is claimed here.
\item No additional authoritative citation or full reference is supplied for the moduli-action criterion or examples.
\item Research attempt 1: result ``reduction,'' problem status at run ``unknown,'' registry status ``partially\_solved,'' rigor label ``proved,'' novelty confidence ``low,'' suggested difficulty L4. The machine-generated aid is not peer reviewed; independently verify before citation.
\end{itemize}
\end{document}
